\documentclass[11pt,a4paper]{article}

% ---------------------------------------------------------
% PACKAGES
% ---------------------------------------------------------
\usepackage[utf8]{inputenc}
\usepackage[T1]{fontenc}
\usepackage{lmodern}
\usepackage{amsmath, amssymb, amsfonts}
\usepackage{bm}
\usepackage{graphicx}
\usepackage{hyperref}
\usepackage{authblk}
\usepackage{geometry}
\usepackage{cite}
\usepackage{physics}
\usepackage{tensor}
\usepackage{enumitem}

\geometry{margin=1in}

% ---------------------------------------------------------
% TITLE
% ---------------------------------------------------------
\title{\textbf{Companion XV: The Early-Universe Gravitational-Wave Spectrum in the Relaxation-Driven Cyclic Cosmology (RDCC)}\\[2mm]
\textbf{Relaxation-Induced Tensor Damping, Sectoral Asymmetry, and CPT-Sector GW Signatures}}

\author[1]{Michael Lehmann}
\affil[1]{\small Independent Researcher, Relaxation-Driven Cyclic Cosmology (RDCC) Framework\\
\texttt{mi.lehmann@gmx.de}}

\date{\small \today}

% ---------------------------------------------------------
% DOCUMENT START
% ---------------------------------------------------------
\begin{document}

\maketitle

\begin{abstract}
The early-Universe gravitational-wave (GW) spectrum provides a uniquely sensitive probe of
the microphysical processes that shaped the primordial plasma. In the Relaxation-Driven Cyclic
Cosmology (RDCC), the universal relaxation mechanism introduced in Companion~X modifies
the evolution of tensor perturbations during the earliest stages of cosmic history. The relaxation
rate $\Gamma(a) = \alpha(a) H(a)$ introduces an additional damping term in the tensor-mode
equation of motion, suppressing the amplitude of primordial gravitational waves on scales that
re-entered the horizon during the relaxation-dominated epoch.

Combined with the sectoral temperature asymmetry between the two CPT-conjugate sectors,
this leads to a characteristic sectoral GW asymmetry and a relaxation-imprinted suppression
feature in the stochastic gravitational-wave background (SGWB). The colder CPT sector
experiences weaker damping and produces a distinct tensor spectrum, while the visible sector
exhibits a relaxation-induced suppression at high frequencies.

In this Companion we derive the relaxation-corrected tensor-mode evolution, compute the
resulting GW transfer function, and obtain the full early-Universe GW spectrum in RDCC. We
incorporate the modernized sectoral thermodynamics from Companion~XIII, the CPT-sector PBH
interpretation from Companion~XIV, and the tensor-modulation structure from Companion~II.
The results establish the SGWB as a central observational probe of the relaxation-driven
dynamics that define the RDCC framework.
\end{abstract}

% ========================================================
% 1 MOTIVATION
% ========================================================

\section{Motivation}

The stochastic gravitational-wave background (SGWB) generated in the early Universe provides
a uniquely sensitive probe of high-energy physics. In standard cosmology, the SGWB is
determined by the primordial tensor spectrum and its subsequent evolution through radiation
and matter domination. The resulting spectrum is smooth, nearly scale-invariant, and free of
dynamical features except for small steps associated with changes in the relativistic degrees of
freedom.

In the Relaxation-Driven Cyclic Cosmology (RDCC), the situation is fundamentally different.
The universal relaxation mechanism introduced in Companion~X modifies the evolution of
super-horizon tensor modes during the earliest stages of cosmic history. When the relaxation
rate satisfies
\begin{equation}
\Gamma(a) = \alpha(a) H(a) \gtrsim H(a),
\end{equation}
tensor modes experience an additional damping term that suppresses their amplitude before
horizon re-entry. This effect leaves a characteristic imprint on the SGWB: a smooth, power-law
suppression at frequencies corresponding to modes that re-entered during the relaxation-dominated
epoch.

A second key ingredient is the sectoral temperature asymmetry between the two CPT-conjugate
sectors. As shown in Companion~X, the temperature ratio
\begin{equation}
\frac{T_{B}}{T_{A}} \simeq 0.5,
\end{equation}
implies different relaxation histories, horizon masses, and damping scales in the two sectors.
Tensor modes in the colder CPT sector experience weaker damping, leading to a sectoral GW
asymmetry.

This Companion develops the early-Universe gravitational-wave sector of RDCC, incorporating
the modernized sectoral thermodynamics from Companion~XIII, the CPT-sector PBH
interpretation from Companion~XIV, and the tensor-modulation structure from Companion~II.
The SGWB emerges as a central observational probe of the relaxation-driven dynamics that
define the RDCC framework.

The structure of this Companion is as follows. Section~2 reviews tensor modes in standard
cosmology. Section~3 derives the relaxation-induced tensor damping in RDCC. Section~4
constructs the relaxation-corrected transfer function. Section~5 presents the full early-Universe
GW spectrum. Section~6 develops the observational predictions, and Section~7 provides the
conclusion.

% ========================================================
% 2 TENSOR MODES IN STANDARD COSMOLOGY
% ========================================================

\section{Tensor Modes in Standard Cosmology}

In standard cosmology, primordial tensor perturbations obey the equation
\begin{equation}
\ddot{h}_{k} + 3H \dot{h}_{k} + \frac{k^{2}}{a^{2}} h_{k} = 0,
\end{equation}
with two distinct dynamical regimes:

\begin{itemize}
\item \textbf{Super-horizon regime} (\(k < aH\)):  
Tensor modes freeze out with constant amplitude,  
\(\dot{h}_{k} \approx 0\).

\item \textbf{Sub-horizon regime} (\(k > aH\)):  
Tensor modes oscillate with physical frequency \(k/a\) and decay as \(a^{-1}\) due to Hubble friction.
\end{itemize}

The primordial tensor spectrum is commonly parameterized as
\begin{equation}
\mathcal{P}_{h}(k) = A_{t} \left( \frac{k}{k_{*}} \right)^{n_{t}},
\end{equation}
where \(A_{t}\) is the tensor amplitude, \(n_{t}\) is the tensor tilt, and \(k_{*}\) is a pivot scale.

The present-day gravitational-wave energy density per logarithmic interval in wavenumber is
\begin{equation}
\Omega_{\mathrm{GW}}(k)
=
\frac{1}{12}
\left( \frac{k}{a_{0} H_{0}} \right)^{2}
\mathcal{P}_{h}(k)\, T^{2}(k),
\end{equation}
where \(T(k)\) is the tensor transfer function.

In standard cosmology, the transfer function is determined by:

\begin{itemize}
\item horizon re-entry during radiation or matter domination,
\item changes in the relativistic degrees of freedom,
\item the transition from radiation to matter domination.
\end{itemize}

Crucially, \textbf{standard cosmology contains no additional damping beyond Hubble friction}.  
The SGWB is therefore smooth and featureless, making it exceptionally sensitive to any new
early-Universe physics.

This provides the baseline against which the relaxation-induced tensor damping of RDCC must
be compared, developed in Section~3.

% ========================================================
% 3 RELAXATION-INDUCED TENSOR DAMPING
% ========================================================

\section{Relaxation-Induced Tensor Damping}

In the Relaxation-Driven Cyclic Cosmology (RDCC), the evolution of tensor perturbations is
modified by the universal relaxation mechanism introduced in Companion~X. The relaxation
rate
\begin{equation}
\Gamma(a) = \alpha(a) H(a)
\end{equation}
acts as an additional friction term in the tensor-mode equation of motion. This modifies the
dynamics of super-horizon tensor modes during the relaxation-dominated epoch and imprints a
characteristic suppression feature on the stochastic gravitational-wave background (SGWB).

The tensor-mode equation in RDCC becomes
\begin{equation}
\ddot{h}_{k} + \left( 3H + \Gamma \right) \dot{h}_{k} + \frac{k^{2}}{a^{2}} h_{k} = 0,
\end{equation}
which can be written as
\begin{equation}
\ddot{h}_{k} + (3 + \alpha)H \dot{h}_{k} + \frac{k^{2}}{a^{2}} h_{k} = 0.
\end{equation}

% --------------------------------------------------------
% 3.1 Super-Horizon Evolution
% --------------------------------------------------------

\subsection{Super-Horizon Evolution}

For modes outside the horizon (\(k < aH\)), the gradient term is negligible and the equation
reduces to
\begin{equation}
\ddot{h}_{k} + (3 + \alpha)H \dot{h}_{k} \approx 0.
\end{equation}

In standard cosmology, tensor modes freeze out with constant amplitude. In RDCC, the
additional relaxation term suppresses the amplitude according to
\begin{equation}
h_{k}(a) \propto a^{-\alpha/2},
\end{equation}
implying that modes which remain super-horizon during the relaxation-dominated epoch
experience a smooth, power-law suppression.

This suppression is the tensor analogue of the relaxation-induced damping of scalar
perturbations derived in Companion~X.

% --------------------------------------------------------
% 3.2 Sub-Horizon Evolution
% --------------------------------------------------------

\subsection{Sub-Horizon Evolution}

Once a mode re-enters the horizon (\(k > aH\)), the oscillatory term dominates and the solution
takes the form
\begin{equation}
h_{k}(a) \propto \frac{1}{a} \exp\!\left[ -\frac{1}{2}\int \alpha(a)\, d\ln a \right]
\cos\!\left( k\int \frac{da}{a^{2}H} \right).
\end{equation}

The usual \(a^{-1}\) decay from Hubble friction is now multiplied by an additional exponential
suppression factor sourced by the relaxation dynamics. This factor depends on the integrated
relaxation history of the mode.

% --------------------------------------------------------
% 3.3 Relaxation-Dominated Epoch
% --------------------------------------------------------

\subsection{Relaxation-Dominated Epoch}

The relaxation-dominated regime is defined by
\begin{equation}
\Gamma(a) \gtrsim H(a),
\qquad\text{or equivalently}\qquad
\alpha(a) \gtrsim 1.
\end{equation}

Modes that re-enter the horizon during this epoch experience the strongest suppression. The
corresponding wavenumbers satisfy
\begin{equation}
k_{\mathrm{rel}} \sim a_{\mathrm{rel}} H_{\mathrm{rel}},
\end{equation}
where \(a_{\mathrm{rel}}\) denotes the scale factor at which \(\alpha(a)\) transitions below unity.

Modes with \(k > k_{\mathrm{rel}}\) re-enter during the relaxation-dominated epoch and are therefore
suppressed, while modes with \(k < k_{\mathrm{rel}}\) re-enter later and retain their standard amplitude.

This produces a smooth, power-law suppression feature in the SGWB.

% --------------------------------------------------------
% 3.4 Sectoral Tensor Asymmetry
% --------------------------------------------------------

\subsection{Sectoral Tensor Asymmetry}

The sectoral temperature asymmetry derived in Companion~X,
\begin{equation}
\frac{T_{B}}{T_{A}} \simeq 0.5,
\end{equation}
implies different relaxation histories in the two sectors:
\begin{equation}
\alpha_{B}(a) < \alpha_{A}(a).
\end{equation}

As a result:

\begin{itemize}
\item tensor modes in the visible sector experience stronger damping,
\item tensor modes in the CPT sector experience weaker damping,
\item the SGWB exhibits a sectoral asymmetry in amplitude and spectral shape.
\end{itemize}

This sectoral tensor asymmetry is a distinctive and testable prediction of RDCC and forms the
basis for the relaxation-corrected transfer function developed in Section~4.

% ========================================================
% 4 RELAXATION-CORRECTED TRANSFER FUNCTION
% ========================================================

\section{Relaxation-Corrected Transfer Function}

The transfer function encodes the evolution of tensor modes from horizon re-entry to the
present day. In standard cosmology, the transfer function is determined by the expansion
history and the changes in the relativistic degrees of freedom. In the Relaxation-Driven Cyclic
Cosmology (RDCC), the universal relaxation mechanism introduces an additional suppression
factor that modifies the transfer function at high frequencies.

The relaxation-corrected transfer function can be written as
\begin{equation}
T^{2}(k)
=
T_{\mathrm{std}}^{2}(k)\,
T_{\mathrm{rel}}^{2}(k),
\end{equation}
where \(T_{\mathrm{std}}(k)\) is the standard transfer function and \(T_{\mathrm{rel}}(k)\) captures the
relaxation-induced suppression.

% --------------------------------------------------------
% 4.1 Standard Transfer Function
% --------------------------------------------------------

\subsection{Standard Transfer Function}

The standard transfer function is given by
\begin{equation}
T_{\mathrm{std}}^{2}(k)
=
\Omega_{r,0}
\left[
1 + \frac{4}{3}\left( \frac{k}{k_{\mathrm{eq}}} \right)
+ \frac{5}{2}\left( \frac{k}{k_{\mathrm{eq}}} \right)^{2}
\right]^{-1},
\end{equation}
where \(k_{\mathrm{eq}}\) is the wavenumber corresponding to matter–radiation equality and
\(\Omega_{r,0}\) is the present-day radiation density fraction.

This expression captures the standard suppression of tensor modes that re-enter during matter
domination.

% --------------------------------------------------------
% 4.2 Relaxation-Induced Suppression Factor
% --------------------------------------------------------

\subsection{Relaxation-Induced Suppression Factor}

The relaxation-induced suppression factor arises from the additional damping term in the
tensor-mode equation of motion. For a mode with wavenumber \(k\), the suppression is given by
\begin{equation}
T_{\mathrm{rel}}(k)
=
\exp\!\left[
-\frac{1}{2}
\int_{\ln a_{\mathrm{hc}}}^{\ln a_{\mathrm{end}}}
\alpha(a)\, d\ln a
\right],
\end{equation}
where \(a_{\mathrm{hc}}\) is the scale factor at horizon crossing and \(a_{\mathrm{end}}\) marks the end of the
relaxation-dominated epoch.

Modes that re-enter the horizon during the relaxation-dominated epoch experience the largest
suppression. Modes that re-enter later are unaffected.

% --------------------------------------------------------
% 4.3 Characteristic Relaxation Scale
% --------------------------------------------------------

\subsection{Characteristic Relaxation Scale}

The characteristic relaxation scale is defined by the wavenumber
\begin{equation}
k_{\mathrm{rel}} = a_{\mathrm{rel}} H_{\mathrm{rel}},
\end{equation}
where \(a_{\mathrm{rel}}\) is the scale factor at which \(\alpha(a)\) transitions below unity. Modes with
\(k > k_{\mathrm{rel}}\) re-enter during the relaxation-dominated epoch and are therefore suppressed,
while modes with \(k < k_{\mathrm{rel}}\) retain their standard amplitude.

The relaxation-induced suppression factor can be approximated as
\begin{equation}
T_{\mathrm{rel}}(k)
\simeq
\left( \frac{k}{k_{\mathrm{rel}}} \right)^{-\alpha_{\mathrm{eff}}},
\qquad
k > k_{\mathrm{rel}},
\end{equation}
where \(\alpha_{\mathrm{eff}}\) is an effective relaxation coefficient that captures the integrated
relaxation history.

% --------------------------------------------------------
% 4.4 Sectoral Transfer Functions
% --------------------------------------------------------

\subsection{Sectoral Transfer Functions}

The sectoral temperature asymmetry derived in Companion~X,
\begin{equation}
\frac{T_{B}}{T_{A}} \simeq 0.5,
\end{equation}
implies different relaxation histories in the two sectors. The relaxation coefficients satisfy
\begin{equation}
\alpha_{B}(a) < \alpha_{A}(a),
\end{equation}
leading to sector-dependent transfer functions:
\begin{equation}
T_{\mathrm{rel}}^{(B)}(k) > T_{\mathrm{rel}}^{(A)}(k).
\end{equation}

The visible sector therefore exhibits a stronger high-frequency suppression, while the CPT
sector retains a larger tensor amplitude.

This sectoral asymmetry is a distinctive and testable prediction of RDCC and forms the basis
for the full early-Universe GW spectrum developed in Section~5.

% ========================================================
% 5 FULL EARLY-UNIVERSE GW SPECTRUM
% ========================================================

\section{Full Early-Universe GW Spectrum}

The relaxation-corrected transfer function derived in Section~4 allows us to construct the full
early-Universe gravitational-wave (GW) spectrum in the Relaxation-Driven Cyclic Cosmology
(RDCC). The present-day GW energy density per logarithmic interval in wavenumber is given by
\begin{equation}
\Omega_{\mathrm{GW}}(k)
=
\frac{1}{12}
\left( \frac{k}{a_{0} H_{0}} \right)^{2}
\mathcal{P}_{h}(k)\,
T_{\mathrm{std}}^{2}(k)\,
T_{\mathrm{rel}}^{2}(k),
\end{equation}
where the relaxation-induced suppression factor $T_{\mathrm{rel}}(k)$ encodes the integrated
relaxation history of each mode.

The resulting spectrum exhibits three characteristic regimes.

% --------------------------------------------------------
% 5.1 Low-Frequency Regime
% --------------------------------------------------------

\subsection{Low-Frequency Regime}

Modes with $k < k_{\mathrm{rel}}$ re-enter the horizon after the relaxation-dominated epoch. Their
amplitude is unaffected by the relaxation mechanism, and the spectrum reduces to the standard
result:
\begin{equation}
\Omega_{\mathrm{GW}}(k)
\simeq
\Omega_{\mathrm{GW}}^{\mathrm{std}}(k),
\qquad
k < k_{\mathrm{rel}}.
\end{equation}

This regime is identical in both sectors and provides a baseline for comparison.

% --------------------------------------------------------
% 5.2 Intermediate Regime
% --------------------------------------------------------

\subsection{Intermediate Regime}

Modes with $k \sim k_{\mathrm{rel}}$ re-enter the horizon near the end of the relaxation-dominated
epoch. These modes experience partial suppression, leading to a smooth transition between the
standard and relaxation-dominated regimes. The spectrum takes the form
\begin{equation}
\Omega_{\mathrm{GW}}(k)
\simeq
\Omega_{\mathrm{GW}}^{\mathrm{std}}(k)
\left( \frac{k}{k_{\mathrm{rel}}} \right)^{-2\alpha_{\mathrm{eff}}},
\qquad
k \sim k_{\mathrm{rel}}.
\end{equation}

This regime encodes the detailed relaxation history and is sensitive to the shape of $\alpha(a)$.

% --------------------------------------------------------
% 5.3 High-Frequency Regime
% --------------------------------------------------------

\subsection{High-Frequency Regime}

Modes with $k > k_{\mathrm{rel}}$ re-enter during the relaxation-dominated epoch and experience
the strongest suppression. The spectrum becomes
\begin{equation}
\Omega_{\mathrm{GW}}(k)
\simeq
\Omega_{\mathrm{GW}}^{\mathrm{std}}(k)
\left( \frac{k}{k_{\mathrm{rel}}} \right)^{-2\alpha_{\mathrm{eff}}},
\qquad
k > k_{\mathrm{rel}}.
\end{equation}

This power-law suppression is the defining signature of relaxation-induced tensor damping.

% --------------------------------------------------------
% 5.4 Sectoral GW Spectra
% --------------------------------------------------------

\subsection{Sectoral GW Spectra}

The sectoral temperature asymmetry derived in Companion~X implies different relaxation
histories in the two sectors:
\begin{equation}
\alpha_{B}(a) < \alpha_{A}(a).
\end{equation}

As a result:

\begin{itemize}
\item the visible sector exhibits a stronger high-frequency suppression,
\item the CPT sector retains a larger tensor amplitude,
\item the SGWB displays a sectoral asymmetry in amplitude and spectral slope.
\end{itemize}

The full sectoral spectra are therefore
\begin{equation}
\Omega_{\mathrm{GW}}^{(A)}(k)
=
\Omega_{\mathrm{GW}}^{\mathrm{std}}(k)
\left( \frac{k}{k_{\mathrm{rel}}^{(A)}} \right)^{-2\alpha_{\mathrm{eff}}^{(A)}},
\end{equation}
\begin{equation}
\Omega_{\mathrm{GW}}^{(B)}(k)
=
\Omega_{\mathrm{GW}}^{\mathrm{std}}(k)
\left( \frac{k}{k_{\mathrm{rel}}^{(B)}} \right)^{-2\alpha_{\mathrm{eff}}^{(B)}},
\end{equation}
with
\begin{equation}
\alpha_{\mathrm{eff}}^{(B)} < \alpha_{\mathrm{eff}}^{(A)},
\qquad
k_{\mathrm{rel}}^{(B)} < k_{\mathrm{rel}}^{(A)}.
\end{equation}

These expressions define the complete early-Universe GW spectrum in RDCC and form the
basis for the observational predictions developed in Section~6.

% ========================================================
% 6 OBSERVATIONAL PREDICTIONS
% ========================================================

\section{Observational Predictions}

The relaxation-induced tensor damping and the sectoral temperature asymmetry of the
Relaxation-Driven Cyclic Cosmology (RDCC) lead to a number of distinctive observational
signatures in the stochastic gravitational-wave background (SGWB). These signatures provide a
direct probe of the relaxation dynamics, the bipartite CPT structure, and the early-Universe
thermodynamics developed in Companions~X, XIII, and XIV.

% --------------------------------------------------------
% 6.1 High-Frequency Suppression
% --------------------------------------------------------

\subsection{High-Frequency Suppression}

The defining prediction of RDCC is a smooth, power-law suppression of the SGWB at high
frequencies:
\begin{equation}
\Omega_{\mathrm{GW}}(k)
\propto
\left( \frac{k}{k_{\mathrm{rel}}} \right)^{-2\alpha_{\mathrm{eff}}},
\qquad
k > k_{\mathrm{rel}}.
\end{equation}

This suppression arises from the relaxation-induced damping of tensor modes that re-enter the
horizon during the relaxation-dominated epoch. The scale \(k_{\mathrm{rel}}\) marks the transition
between the standard and relaxation-modified regimes.

% --------------------------------------------------------
% 6.2 Sectoral GW Asymmetry
% --------------------------------------------------------

\subsection{Sectoral GW Asymmetry}

The sectoral temperature ratio
\begin{equation}
\frac{T_{B}}{T_{A}} \simeq 0.5
\end{equation}
implies different relaxation histories in the two sectors. As a result:

\begin{itemize}
\item the visible sector exhibits a stronger high-frequency suppression,
\item the CPT sector retains a larger tensor amplitude,
\item the SGWB displays a sector-dependent spectral slope.
\end{itemize}

This sectoral asymmetry is a distinctive and testable prediction of RDCC.

% --------------------------------------------------------
% 6.3 Connection to PBH Physics
% --------------------------------------------------------

\subsection{Connection to PBH Physics}

The same relaxation dynamics that suppress tensor modes also suppress PBH formation in the
visible sector and enhance PBH formation in the CPT sector. The SGWB and the PBH sector
therefore provide complementary probes of the relaxation-driven early-Universe physics.

The characteristic scales \(k_{\mathrm{rel}}^{(A/B)}\) are directly related to the PBH mass scales
derived in Companion~XIV.

% --------------------------------------------------------
% 6.4 Experimental Prospects
% --------------------------------------------------------

\subsection{Experimental Prospects}

Future gravitational-wave observatories will be able to probe the relaxation-induced features of
the SGWB:

\begin{itemize}
\item \textbf{LISA}: sensitive to the intermediate regime near \(k_{\mathrm{rel}}\),
\item \textbf{DECIGO/BBO}: sensitive to the high-frequency suppression,
\item \textbf{Einstein Telescope / Cosmic Explorer}: complementary constraints at lower frequencies.
\end{itemize}

A detection of a smooth, power-law suppression at high frequencies would provide strong
evidence for relaxation-driven tensor damping.

% --------------------------------------------------------
% 6.5 Consistency with RDCC Thermodynamics
% --------------------------------------------------------

\subsection{Consistency with RDCC Thermodynamics}

The SGWB predictions are fully consistent with the sectoral thermodynamics developed in
Companion~XIII and the PBH sector developed in Companion~XIV. The relaxation-induced
tensor suppression, the sectoral asymmetry, and the PBH mass hierarchy all arise from the same
underlying relaxation dynamics and the bipartite CPT structure of RDCC.

% ========================================================
% 7 CONCLUSION
% ========================================================

\section{Conclusion}

In this Companion we have developed the early-Universe gravitational-wave (GW) sector of the
Relaxation-Driven Cyclic Cosmology (RDCC). The universal relaxation mechanism introduced in
Companion~X modifies the evolution of tensor perturbations during the earliest stages of cosmic
history, introducing an additional damping term that suppresses the amplitude of primordial
gravitational waves on scales that re-entered the horizon during the relaxation-dominated epoch.

The resulting stochastic gravitational-wave background (SGWB) exhibits a smooth, power-law
suppression at high frequencies, with the characteristic scale set by the relaxation transition
\(k_{\mathrm{rel}}\). This suppression is the tensor analogue of the relaxation-induced damping of scalar
perturbations and provides a direct probe of the relaxation dynamics that define the RDCC
framework.

The sectoral temperature asymmetry between the two CPT-conjugate sectors leads to different
relaxation histories and therefore to a sectoral GW asymmetry. Tensor modes in the visible
sector experience stronger damping, while modes in the CPT sector retain a larger amplitude.
This sectoral asymmetry is a distinctive and testable prediction of RDCC.

The full early-Universe GW spectrum constructed in this Companion incorporates the
modernized sectoral thermodynamics from Companion~XIII, the CPT-sector PBH interpretation
from Companion~XIV, and the tensor-modulation structure from Companion~II. Together, these
results establish the SGWB as a central observational probe of the relaxation-driven early-Universe
physics that underlies the RDCC framework.

Future gravitational-wave observatories such as LISA, DECIGO, BBO, and the Einstein
Telescope will be able to test the relaxation-induced suppression feature and the sectoral GW
asymmetry. A detection of these signatures would provide strong evidence for the relaxation
mechanism and the bipartite CPT structure that define RDCC.


\end{document}
